\subsection{Robustness, Shift, and Distributional Uncertainty}

Learning systems are deployed under distributions that may differ from the training law.

\begin{definition} \label{def:sl-distribution-shift-types} \hlt{Distribution Shift}\\
Common forms of distribution shift include:
\begin{enumerate}[label=(\roman*)]
\setlength{\itemsep}{0pt}
\item Covariate shift changes $P_X$ while preserving $P_{Y\mid X}$.
\item Label shift changes $P_Y$ while preserving $P_{X\mid Y}$.
\item Concept shift changes the relationship between inputs and outputs.
\end{enumerate}
\end{definition}

\begin{definition} \label{def:sl-adversarial-example} \hlt{Adversarial Example}\\
An adversarial example is a perturbed input $x'$ near $x$ for which the model output changes in a harmful way.  Adversarial training adds such perturbed inputs to the training objective to improve local robustness.
\end{definition}

\begin{definition} \label{def:sl-distributionally-robust-optimization} \hlt{Distributionally Robust Optimization}\\
Distributionally robust optimization minimizes worst-case risk over an uncertainty set $\mathcal{U}$ of plausible distributions:
\[
\inf_f\sup_{Q\in\mathcal{U}}\mathbb{E}_{Q}\ell(f(X),Y).
\]
\end{definition}

\begin{remark} \label{rem:sl-uncertainty-calibration} \hlt{Uncertainty}\\
Reliable systems distinguish aleatoric noise, epistemic uncertainty, and out-of-distribution uncertainty.  Calibration and abstention policies are part of robustness, not decorative outputs.
\end{remark}
